\documentclass[11pt]{article}
\usepackage{amsmath,amssymb,amsthm}
\newtheorem{definition}{Definition}
\newtheorem{lemma}{Lemma}
\newtheorem{corollary}{Corollary}
\begin{document}

\begin{definition}[Explicit half-edge labeling]
For each base edge $e=\{u,v\}\in E(B)$, introduce half-edge labels
\[
h_{u,e},h_{v,e}\in\mathbb F_2.
\]
For each base vertex $u$, define the local CFI state
\[
\sigma_u:E(u)\to\mathbb F_2
\]
subject to the parity constraint
\[
\sum_{e\ni u}\sigma_u(e)=0.
\]
\end{definition}

\begin{definition}[Transport cocycle from matching constraint]
For adjacent $u\sim v$, define
\[
\tau_{uv}:=\sigma_u(\{u,v\})+\sigma_v(\{u,v\})\in\mathbb F_2.
\]
\end{definition}

\begin{lemma}[Cocycle identity]
\[
\tau_{uv}+\tau_{vw}+\tau_{wu}=0.
\]
This follows by cancellation of shared half-edge contributions on triple overlaps.
\end{lemma}

\begin{lemma}[Kernel triviality]
\[
\ker(\Phi_{\mathrm{explicit}})=0.
\]
This follows by normalization of local states modulo transport.
\end{lemma}

\begin{definition}[Inverse map]
Define
\[
\Psi:H^1(B;\mathbb F_2)\to \mathcal H_R(G),
\qquad
\Psi([\beta])=[(\sigma_u^\beta)_u],
\]
where
\[
\sigma_u^\beta(e):=\beta(e)\qquad (e\ni u),
\]
followed by parity normalization inside each CFI gadget.
\end{definition}

\begin{lemma}
\[
\Phi_{\mathrm{explicit}}\circ\Psi=\mathrm{id}.
\]
\end{lemma}

\begin{lemma}
\[
\Psi\circ\Phi_{\mathrm{explicit}}=\mathrm{id}.
\]
\end{lemma}

\begin{corollary}
\[
\mathcal H_R(G)\cong H^1(B;\mathbb F_2).
\]
\end{corollary}

\end{document}
